%------------------------------------------------------------------------------%
\section{Integrais Trigonométricas}
\label{sec:Integais_Trigonometricas}

%------------------------------------------------------------------------------%
Fórmulas úteis de trigonometria
\begin{enumerate}
  \item
    Relação fundamental
    \[
      \sin^2(x)+ \cos^2(x)=1
    \]

  \item
    Decorrências da relação fundamental
    \begin{gather*}
      \tan^2(x)+ 1  = \sec^2(x) \\
      1 + \cot^2(x) = \csc^2(x)
    \end{gather*}

  \item
    Arco duplo
    \begin{gather*}
      \sin(2x) = 2\sin(x) \cos(x) \\
      \cos(2x) = \cos^2(x)-\sin^2(x)
    \end{gather*}
  \item
    Arco metade
    \begin{gather*}
      \sin^2(x)= \frac{1}{2}(1-\cos(2x)) \\
      \cos^2(x)= \frac{1}{2}\big(1+\cos(2x)\big)
    \end{gather*}
\end{enumerate}
%------------------------------------------------------------------------------%

Estratégias para calcular \(\int \sin^m(x) \cos^n(x)dx\)
\begin{enumerate}
  \item Se a potência de $\sin(x)$ for ímpar:
    \begin{enumerate}
      \item Guarda um fator $\sin(x)$ e use que $\sin^2(x)=1- \cos^2(x)$.
      \item Faça a substituição $u=\cos(x)$.
    \end{enumerate}

  \item Se a potência de $\cos(x)$ for ímpar:
    \begin{enumerate}
      \item Guarda um fator $\cos(x)$ e use que $\cos^2(x)=1- \sin^2(x)$.
      \item Faça a substituição $u=\sin(x)$.
    \end{enumerate}

  \item Se as potências de seno e cosseno forem pares, use uma das transformações:
    \[
      \sin^2(x)= \frac{1}{2}\big(1-\cos(2x)\big) \qquad\text{ou}\qquad
      \cos^2(x)= \frac{1}{2}\big(1+\cos(2x)\big)
    \]
\end{enumerate}
Algumas vezes a transformação $\sin(2x)=2\sin(x)\cos(x)$ pode ser útil.
%------------------------------------------------------------------------------%

%------------------------------------------------------------------------------%
Estratégias para calcular \(\int \tan^m(x) \sec^n(x)dx\):
\begin{enumerate}
  \item Se a potência da $\sec(x)$ é par ($n=2k$, $k \geq 2$).
        \begin{enumerate}
          \item Guarda um fator $\sec^2(x)$ e use que $\sec^2(x)=1+\tan^2(x)$.
          \item Faça a substituição $u=\tan(x)$.
        \end{enumerate}

  \item Se a potência de $\tan(x)$ for ímpar:
        \begin{enumerate}
          \item Guarda um fator $\sec(x)\tan(x)$ e use que $\tan^2(x)=\sec^2(x)-1$.
          \item Faça a substituição $u=\sec(x)$.
        \end{enumerate}
        Note que a mesma ideia pode ser usada para integrais do tipo
        \(\int \cot^m(x) \csc^m(x)dx\)

  \item Os outros casos não tem uma regra geral e requer manipulações
        trigonométricas específicas para cada caso.
\end{enumerate}
%------------------------------------------------------------------------------%

%------------------------------------------------------------------------------%
\begin{example}
  Calcule \(\int \sin^3(x)\cos^2(x)dx\)
  \solution
  Aqui, temos uma integral do tipo
  $\int \sin^m(x) \cos^n(x)dx$ e a potência de $\sin(x)$ é ímpar,
  por isso, guardamos um fator $\sin(x)$ e usamos a transformação
  $\sin^2(x)= 1- \cos^2(x)$, decorrente da relação fundamental. Assim,
  \[
    \sin^3(x)\cos^2(x)= \sin^2(x)\cos^2(x)\sin(x)= (1-\cos^2)cos^2(x)\sin(x)
  \]
  e podemos reescrever a integral
  \[
    \int \sin^3(x)\cos^2(x)dx=\int  (1-\cos^2)cos^2(x)\sin(x) dx
  \]
  Usando a substituição simples $u= \cos(x)$, temos que $du=-\sin(x)dx$ e, assim,
  \begin{align*}
    \int \sin^3(x)\cos^2(x)dx
    & = \int  (1-\cos^2(x))cos^2(x)\sin(x) dx \\
    & = \int (1-u^2)u^2(-du)=- \int u^2-u^4 du \\
    & = \frac{u^3}{3}- \frac{u^5}{5}+C \\
    & = \frac{\cos^3(x)}{3}- \frac{\cos^5(x)}{5}+C
  \end{align*}
\end{example}
%------------------------------------------------------------------------------%

Vamos pensar um pouco sobre o porque que esse procedimento que fizemos é adequado.
Você poderia questionar: Por que não usamos a transformação
$\sin^2(x)=1-\cos^2(x)$?
Bom, caso fizéssimos essa mudança, teríamos a seguinte integral resultante
\[
  \int \sin^3(x)\cos^2(x)dx
  = \int   \sin^3(x)(1-\sin^2(x))dx
  = \int [ \sin^3(x)- \sin^5(x)]dx
\]
ou seja, teríamos uma expressão em termos de $\sin(x)$ sem nenhum fator
extra $\cos(x)$ necessário para usarmos substituição simples.

%------------------------------------------------------------------------------%
\begin{example}
  Encontre \(\int \cos^3(x)dx\)
  \solution
  Aqui temos uma integral que envolve apenas potência de
  $\cos(x)$, mas mesmo assim usamos o mesmo procedimento.
  Como o expoente do $\cos(x)$ é ímpar, então guardamos um fator
  $\cos(x)$ e usamos a transformação $\cos^2(x)= 1-\sin^2(x)$ e assim,
  usamos a substituição simples $u= \sin(x) dx \Leftrightarrow du=\cos(x)dx$.
  Com isso,
  \begin{align*}
    \int \cos^3(x)dx & = \int \cos^2(x) \cos(x)dx       \\
                     & = \int (1-\sin^2(x))\cos(x)dx    \\
                     & = \int (1-u^2)du                 \\
                     & = u- \frac{u^3}{3}+C             \\
                     & = \cos(x)- \frac{\cos^3(x)}{3}+C
  \end{align*}
\end{example}
%------------------------------------------------------------------------------%

%------------------------------------------------------------------------------%
\begin{example}
  Calcule \(\int_0^{\pi} \sin^2(x) dx\)
  \solution
  Aqui temos apenas uma potência par de $\sin(x)$,
  portanto usamos a transformação
  $\sin^2(x)= \frac{1}{2}(1-\cos(2x))$.
  Com isso,
  \begin{align*}
    \int_0^{\pi} \sin^2(x) dx
    & = \int_0^\pi \frac{1}{2}(1-\cos(2x)) dx \\
    & = \frac{1}{2} \int_0^\pi (1-\cos(2x)) dx \\
    & = \frac{1}{2} \left[  x - \frac{\sin(2x)}{2}\right] \evalat{0}{\pi} \\
    & = \frac{1}{2}
        \left[
          \left( \pi - \frac{\sin(2\pi)}{2} \right)
         -\left(   0 - \frac{\sin(0)}{2}    \right)
        \right] \\
    & = \frac{\pi}{2}
  \end{align*}
  onde a penúltima igualdade decorre do Teorema fundamental do cálculo.
\end{example}
%------------------------------------------------------------------------------%

%------------------------------------------------------------------------------%
\begin{example}
  Calcule \(\int \sin^2(x) \cos^2(x) dx\)
  \solution
  Como aqui tanto a potência do $\sin(x)$ quanto a potência
  do $\cos(x)$ são pares, vamos usar as transformações do arco metade.
  Assim,
  \begin{align*}
    \int \sin^2(x) \cos^2(x) dx
    & = \int \frac{1}{2}(1- \cos(2x)) \frac{1}{2}(1+ \cos(2x)) dx \\
    & = \frac{1}{4} \int (1- \cos^2(2x)) dx
  \end{align*}
  Precisamos mais uma vez usar a fórmula do arco metade escrevendo
  $\cos^2(2x)= \frac{1}{2}(1+ \cos(4x))$, portanto:
  \begin{align*}
    \int \sin^2(x) \cos^2(x) dx
    & = \frac{1}{4} \int (1- \cos^2(2x)) dx \\
    & = \frac{1}{4} \int 1-\frac{1}{2}(1+ \cos(4x))dx \\
    & = \frac{1}{4} \int   \frac{1}{2} - \frac{1}{2} \cos(4x) dx \\
    & = \frac{1}{2} \left( \frac{x}{2} - \frac{1}{8} \sin(4x))\right)+ C
  \end{align*}
\end{example}
%------------------------------------------------------------------------------%

%------------------------------------------------------------------------------%
\begin{example}
  Encontre \(\int \tan^6(x) \sec^4(x) dx\)
  \solution
  Nesse exercícios vamos separar um fator $\sec^2(x)$ e
  expressar $\sec^2(x)$ em termos da tangente, usando a identidade
  $\sec^2(x)=1+ \tan^2(x)$.
  O próximo passo é usar substituição
  $u=\tan(x)$, de modo que $du=\sec^2(x)dx$:
  \begin{align*}
    \int \tan^6(x) \sec^4(x) dx
    & = \int \tan^6(x) \sec^2(x)  \sec^2(x)dx  \\
    & = \int \tan^6(x) (1+ \tan^2(x))  \sec^2(x)dx \\
    & = \int u^6(1+u^2) du \\
    & = \frac{u^7}{7}+ \frac{u^9}{9}+ C \\
    & = \frac{\tan^7(x)}{7}+ \frac{\tan^9(x)}{9}+ C
  \end{align*}
\end{example}
%------------------------------------------------------------------------------%

%------------------------------------------------------------------------------%
\begin{example}
  Encontre \(\int \tan^5(x) \sec^7(x) dx\)
  \solution
  Nesse caso, temos que a potência de $\tan(x)$ é ímpar, então guardamos
  o fator $\sec(x)\tan(x)$ e usamos a tranformação
  $\tan^2(x)= \sec^2(x)-1$.
  Para resolvermos a integral resultante usamos a substituição
  $u= \sec(x)$, donde $du= \sec(x)\tan(x)dx$
  \begin{align*}
    \int \tan^5(x) \sec^7(x) dx
    & = \int \tan^4(x) \sec^6(x) \sec(x) \tan(x) dx \\
    & = \big(\tan^2(x)\big)^2 \sec^6(x) \sec(x) \tan(x) dx \\
    & = \int (\sec^2(x)-1)^2\sec^6(x) \sec(x) \tan(x) dx \\
    & = \int (u^2-1)^2 u^6 du \\
    & = \int (u^4-2u^2+1)u^6 du \\
    & = \int (u^{10}-2u^8+u^6 )du \\
    & = \frac{u^{11}}{11}-\frac{2u^9}{9}+ \frac{u^7}{7}+C \\
    & = \frac{\sec^{11}(x)}{11}-\frac{2\sec^9(x)}{9}+ \frac{\sec^7(x)}{7}+C
  \end{align*}
\end{example}
%------------------------------------------------------------------------------%

%------------------------------------------------------------------------------%
\begin{example}
  Calcule \(\int \tan^3(x) dx\)
  \solution
  Aqui temos apenas o termo $\tan(x)$, então usamos
  $\tan^2(x)= \sec^2(x)-1$ para reescrever o fator
  $\tan^2(x)$ em termos de $\sec^2(x)$
  \begin{align*}
    \int \tan^3(x) dx
    & = \int \tan(x) \tan^2(x) dx \\
    & = \int \tan(x)\big(\sec^2(x)-1\big) dx \\
    & = \int \tan(x)\sec^2(x) dx- \int \tan(x) dx
  \end{align*}
  Para calcularmos $\int tan(x)\sec^2(x) dx$
  usamos a substituição simples $u=\tan(x)$,
  de modo que $du= \sec^2(x)dx$ e assim,
  $\int tan(x)\sec^2(x) dx= \int udu= \frac{u^2}{2}+ C= \frac{\tan^2(x)}{2}$.
  Logo,
  \begin{align*}
    \int \tan^3(x) dx
    & = \int \tan(x) \tan^2(x) dx \\
    & = \int \tan(x)\big(\sec^2(x)-1\big)dx \\
    & = \int \tan(x)\sec^2(x) dx- \int \tan(x) dx \\
    & = \frac{\tan^2(x)}{2}- \ln\abs{\sec(x)}+C
  \end{align*}
\end{example}
%------------------------------------------------------------------------------%

%------------------------------------------------------------------------------%
